\begin{register}{H}{ASRC\_CHNL0\_CTRL\_CFG\_REG}{0x0000}\label{regdesc:ASRCCHNL0CTRLCFGREG}
\regfield{(reserved)}{22}{10}{0000000000000000000000}%
\regfield{ASRC\_CHNL0\_SEL}{1}{9}{{0}}%
\regfield{ASRC\_CHNL0\_MODE}{2}{7}{{0}}%
\regfield{(reserved)}{1}{6}{0}%
\regfield{ASRC\_CHNL0\_FRAC\_AHEAD}{1}{5}{{0}}%
\regfield{ASRC\_CHNL0\_RS2\_STG0\_MODE}{1}{4}{{0}}%
\regfield{ASRC\_CHNL0\_RS2\_STG1\_MODE}{1}{3}{{0}}%
\regfield{ASRC\_CHNL0\_FRAC\_BYPASS}{1}{2}{{1}}%
\regfield{ASRC\_CHNL0\_RS2\_STG0\_BYPASS}{1}{1}{{1}}%
\regfield{ASRC\_CHNL0\_RS2\_STG1\_BYPASS}{1}{0}{{1}}%
\reglabel{Reset}\regnewline%
\begin{regdesc}\begin{reglist}
\label{fielddesc:ASRCCHNL0RS2STG1BYPASS}\item [ASRC\_CHNL0\_RS2\_STG1\_BYPASS] 设置此位以绕过通道 0 中的第 1 阶重采样器。 (R/W)
\label{fielddesc:ASRCCHNL0RS2STG0BYPASS}\item [ASRC\_CHNL0\_RS2\_STG0\_BYPASS] 设置此位以绕过通道 0 中的第 0 阶重采样器。 (R/W)
\label{fielddesc:ASRCCHNL0FRACBYPASS}\item [ASRC\_CHNL0\_FRAC\_BYPASS] 设置此位以绕过通道 0 中的分数重采样器。 (R/W)
\label{fielddesc:ASRCCHNL0RS2STG1MODE}\item [ASRC\_CHNL0\_RS2\_STG1\_MODE] 写入此位以配置通道 0 中第 1 阶重采样器的模式：\\
0：2 倍插值\\
1：2 倍抽取\\
(R/W)
\label{fielddesc:ASRCCHNL0RS2STG0MODE}\item [ASRC\_CHNL0\_RS2\_STG0\_MODE] 写入此位以配置通道 0 中第 0 阶重采样器的模式：\\
0：2 倍插值\\
1：2 倍抽取\\
(R/W)
\label{fielddesc:ASRCCHNL0FRACAHEAD}\item [ASRC\_CHNL0\_FRAC\_AHEAD] 设置此位以将分数重采样器移至通道 0 中的第 1 阶位置；当输入频率高于输出频率时置 1。 (R/W)
\label{fielddesc:ASRCCHNL0MODE}\item [ASRC\_CHNL0\_MODE] 写入该字段以配置通道模式：\\
0：输入和输出均为单声道\\
1：输入和输出均为双声道\\
2：输入为单声道，输出为双声道\\
3：输入为双声道，输出为单声道\\
(R/W)
\label{fielddesc:ASRCCHNL0SEL}\item [ASRC\_CHNL0\_SEL] 写入此位以选择需要处理的 16 位数据。 (R/W)
\end{reglist}\end{regdesc}
\end{register}


\begin{register}{H}{ASRC\_CHNL0\_FRAC\_COEFF\_REG}{0x0004}\label{regdesc:ASRCCHNL0FRACCOEFFREG}
\regfield{(reserved)}{12}{20}{000000000000}%
\regfield{ASRC\_CHNL0\_FRAC\_L}{10}{10}{{0}}%
\regfield{ASRC\_CHNL0\_FRAC\_M}{10}{0}{{0}}%
\reglabel{Reset}\regnewline%
\begin{regdesc}\begin{reglist}
\label{fielddesc:ASRCCHNL0FRACM}\item [ASRC\_CHNL0\_FRAC\_M] 写入该字段以指定通道 0 中分数重采样因子的分母，与 ASRC\_CHNL0\_FRAC\_L 互质。 (R/W)
\label{fielddesc:ASRCCHNL0FRACL}\item [ASRC\_CHNL0\_FRAC\_L] 写入该字段以指定通道 0 中分数重采样因子的分子，与 ASRC\_CHNL0\_FRAC\_M 互质。 (R/W)
\end{reglist}\end{regdesc}
\end{register}


\begin{register}{H}{ASRC\_CHNL0\_FRAC\_RECIPL\_REG}{0x0008}\label{regdesc:ASRCCHNL0FRACRECIPLREG}
\regfield{(reserved)}{12}{20}{000000000000}%
\regfield{ASRC\_CHNL0\_FRAC\_RECIPL}{20}{0}{{0}}%
\reglabel{Reset}\regnewline%
\begin{regdesc}\begin{reglist}
\label{fielddesc:ASRCCHNL0FRACRECIPL}\item [ASRC\_CHNL0\_FRAC\_RECIPL] 写入该字段以设置通道 0 中的值，计算公式为 \(\left\lfloor \frac{2^{19} + L}{2L} \right\rfloor\)。 (R/W)
\end{reglist}\end{regdesc}
\end{register}


\begin{register}{H}{ASRC\_CHNL0\_RESET\_REG}{0x000C}\label{regdesc:ASRCCHNL0RESETREG}
\regfield{(reserved)}{31}{1}{0000000000000000000000000000000}%
\regfield{ASRC\_CHNL0\_RESET}{1}{0}{{0}}%
\reglabel{Reset}\regnewline%
\begin{regdesc}\begin{reglist}
\label{fielddesc:ASRCCHNL0RESET}\item [ASRC\_CHNL0\_RESET] 设置此位以复位通道 0。 (WT)
\end{reglist}\end{regdesc}
\end{register}


\begin{register}{H}{ASRC\_CHNL0\_START\_REG}{0x0010}\label{regdesc:ASRCCHNL0STARTREG}
\regfield{(reserved)}{31}{1}{0000000000000000000000000000000}%
\regfield{ASRC\_CHNL0\_START}{1}{0}{{0}}%
\reglabel{Reset}\regnewline%
\begin{regdesc}\begin{reglist}
\label{fielddesc:ASRCCHNL0START}\item [ASRC\_CHNL0\_START] 设置此位以启动通道 0。 (R/W)
\end{reglist}\end{regdesc}
\end{register}


\begin{register}{H}{ASRC\_CHNL0\_DIN\_CFG\_REG}{0x0014}\label{regdesc:ASRCCHNL0DINCFGREG}
\regfield{ASRC\_CHNL0\_IN\_LEN}{24}{8}{{0}}%
\regfield{(reserved)}{6}{2}{000000}%
\regfield{ASRC\_CHNL0\_IN\_CNT\_CLR}{1}{1}{{0}}%
\regfield{ASRC\_CHNL0\_IN\_CNT\_ENA}{1}{0}{{0}}%
\reglabel{Reset}\regnewline%
\begin{regdesc}\begin{reglist}
\label{fielddesc:ASRCCHNL0INCNTENA}\item [ASRC\_CHNL0\_IN\_CNT\_ENA] 设置此位以启用输入数据字节计数器。 (R/W)
\label{fielddesc:ASRCCHNL0INCNTCLR}\item [ASRC\_CHNL0\_IN\_CNT\_CLR] 写入此位以清除输入数据字节计数器。 (WT)
\label{fielddesc:ASRCCHNL0INLEN}\item [ASRC\_CHNL0\_IN\_LEN] 写入该字段以指定从 DMA 接收的数据字节数。 (R/W)
\end{reglist}\end{regdesc}
\end{register}


\begin{register}{H}{ASRC\_CHNL0\_DOUT\_CFG\_REG}{0x0018}\label{regdesc:ASRCCHNL0DOUTCFGREG}
\regfield{ASRC\_CHNL0\_OUT\_LEN}{24}{8}{{0}}%
\regfield{(reserved)}{3}{5}{000}%
\regfield{ASRC\_CHNL0\_OUT\_LEN\_COMP}{1}{4}{{0}}%
\regfield{ASRC\_CHNL0\_OUT\_CNT\_CLR}{1}{3}{{0}}%
\regfield{ASRC\_CHNL0\_OUT\_CNT\_ENA}{1}{2}{{0}}%
\regfield{ASRC\_CHNL0\_OUT\_EOF\_GEN\_MODE}{2}{0}{{0}}%
\reglabel{Reset}\regnewline%
\begin{regdesc}\begin{reglist}
\label{fielddesc:ASRCCHNL0OUTEOFGENMODE}\item [ASRC\_CHNL0\_OUT\_EOF\_GEN\_MODE] 写入该字段以指定写入 DMA 的 EOF 类型：\\
0：计数器 EOF\\
1：DMA 输入 EOF\\
2：计数器 EOF 和 DMA 输入 EOF 均写入\\
3：无 EOF 写入\\
(R/W)
\label{fielddesc:ASRCCHNL0OUTCNTENA}\item [ASRC\_CHNL0\_OUT\_CNT\_ENA] 设置此位以启用输出数据字节计数器。 (R/W)
\label{fielddesc:ASRCCHNL0OUTCNTCLR}\item [ASRC\_CHNL0\_OUT\_CNT\_CLR] 写入此位以清除输出数据字节计数器。 (WT)
\label{fielddesc:ASRCCHNL0OUTLENCOMP}\item [ASRC\_CHNL0\_OUT\_LEN\_COMP] 设置此位以启用输出数据字节计数器补偿，当使用分数重采样器并进行 2 倍抽取且 ASRC\_CHNL0\_OUT\_CNT ≥ ASRC\_CHNL0\_OUT\_LEN 时生效。 (R/W)
\label{fielddesc:ASRCCHNL0OUTLEN}\item [ASRC\_CHNL0\_OUT\_LEN] 写入该字段以指定传输到 DMA 的数据字节数，当计数器到达末尾时，EOF（结束标志）将被置位。 (R/W)
\end{reglist}\end{regdesc}
\end{register}


\begin{register}{H}{ASRC\_CHNL0\_FIFO\_CTRL\_REG}{0x001C}\label{regdesc:ASRCCHNL0FIFOCTRLREG}
\regfield{(reserved)}{30}{2}{000000000000000000000000000000}%
\regfield{ASRC\_CHNL0\_OUTFIFO\_RESET}{1}{1}{{0}}%
\regfield{ASRC\_CHNL0\_INFIFO\_RESET}{1}{0}{{0}}%
\reglabel{Reset}\regnewline%
\begin{regdesc}\begin{reglist}
\label{fielddesc:ASRCCHNL0INFIFORESET}\item [ASRC\_CHNL0\_INFIFO\_RESET] 写入此位以复位输入 FIFO。 (WT)
\label{fielddesc:ASRCCHNL0OUTFIFORESET}\item [ASRC\_CHNL0\_OUTFIFO\_RESET] 写入此位以复位输出 FIFO。 (WT)


\end{reglist}\end{regdesc}
\end{register}


\begin{register}{H}{ASRC\_CHNL1\_CTRL\_CFG\_REG}{0x003C}\label{regdesc:ASRCCHNL1CTRLCFGREG}
\regfield{(reserved)}{22}{10}{0000000000000000000000}%
\regfield{ASRC\_CHNL1\_SEL}{1}{9}{{0}}%
\regfield{ASRC\_CHNL1\_MODE}{2}{7}{{0}}%
\regfield{(reserved)}{1}{6}{0}%
\regfield{ASRC\_CHNL1\_FRAC\_AHEAD}{1}{5}{{0}}%
\regfield{ASRC\_CHNL1\_RS2\_STG0\_MODE}{1}{4}{{0}}%
\regfield{ASRC\_CHNL1\_RS2\_STG1\_MODE}{1}{3}{{0}}%
\regfield{ASRC\_CHNL1\_FRAC\_BYPASS}{1}{2}{{1}}%
\regfield{ASRC\_CHNL1\_RS2\_STG0\_BYPASS}{1}{1}{{1}}%
\regfield{ASRC\_CHNL1\_RS2\_STG1\_BYPASS}{1}{0}{{1}}%
\reglabel{Reset}\regnewline%
\begin{regdesc}\begin{reglist}
\label{fielddesc:ASRCCHNL1RS2STG1BYPASS}\item [ASRC\_CHNL1\_RS2\_STG1\_BYPASS] 设置此位以绕过通道 1 中的第 1 阶重采样器。 (R/W)
\label{fielddesc:ASRCCHNL1RS2STG0BYPASS}\item [ASRC\_CHNL1\_RS2\_STG0\_BYPASS] 设置此位以绕过通道 1 中的第 0 阶重采样器。 (R/W)
\label{fielddesc:ASRCCHNL1FRACBYPASS}\item [ASRC\_CHNL1\_FRAC\_BYPASS] 设置此位以绕过通道 1 中的分数重采样器。 (R/W)
\label{fielddesc:ASRCCHNL1RS2STG1MODE}\item [ASRC\_CHNL1\_RS2\_STG1\_MODE] 写入此位以配置通道 1 中第 1 阶重采样器的模式：\\
0：2 倍插值\\
1：2 倍抽取\\
(R/W)
\label{fielddesc:ASRCCHNL1RS2STG0MODE}\item [ASRC\_CHNL1\_RS2\_STG0\_MODE] 写入此位以配置通道 1 中第 0 阶重采样器的模式：\\
0：2 倍插值\\
1：2 倍抽取\\
(R/W)
\label{fielddesc:ASRCCHNL1FRACAHEAD}\item [ASRC\_CHNL1\_FRAC\_AHEAD] 设置此位以将分数重采样器移至通道 1 中的第 1 阶位置；当输入频率高于输出频率时置 1。 (R/W)
\label{fielddesc:ASRCCHNL1MODE}\item [ASRC\_CHNL1\_MODE] 写入该字段以配置通道模式：\\
0：输入和输出均为单声道\\
1：输入和输出均为双声道\\
2：输入为单声道，输出为双声道\\
3：输入为双声道，输出为单声道\\
(R/W)
\label{fielddesc:ASRCCHNL1SEL}\item [ASRC\_CHNL1\_SEL] 写入此位以选择需要处理的 16 位数据。 (R/W)
\end{reglist}\end{regdesc}
\end{register}



\begin{register}{H}{ASRC\_CHNL1\_FRAC\_COEFF\_REG}{0x0040}\label{regdesc:ASRCCHNL1FRACCOEFFREG}
\regfield{(reserved)}{12}{20}{000000000000}%
\regfield{ASRC\_CHNL1\_FRAC\_L}{10}{10}{{0}}%
\regfield{ASRC\_CHNL1\_FRAC\_M}{10}{0}{{0}}%
\reglabel{Reset}\regnewline%
\begin{regdesc}\begin{reglist}
\label{fielddesc:ASRCCHNL1FRACM}\item [ASRC\_CHNL1\_FRAC\_M] 写入该字段以指定通道 1 中分数重采样因子的分母，与 ASRC\_CHNL1\_FRAC\_L 互质。 (R/W)
\label{fielddesc:ASRCCHNL1FRACL}\item [ASRC\_CHNL1\_FRAC\_L] 写入该字段以指定通道 1 中分数重采样因子的分子，与 ASRC\_CHNL1\_FRAC\_M 互质。 (R/W)


\end{reglist}\end{regdesc}
\end{register}


\begin{register}{H}{ASRC\_CHNL1\_FRAC\_RECIPL\_REG}{0x0044}\label{regdesc:ASRCCHNL1FRACRECIPLREG}
\regfield{(reserved)}{12}{20}{000000000000}%
\regfield{ASRC\_CHNL1\_FRAC\_RECIPL}{20}{0}{{0}}%
\reglabel{Reset}\regnewline%
\begin{regdesc}\begin{reglist}
\label{fielddesc:ASRCCHNL1FRACRECIPL}\item [ASRC\_CHNL1\_FRAC\_RECIPL] 写入该字段以设置通道 1 中的值，计算公式为  \(\left\lfloor \frac{2^{19} + L}{2L} \right\rfloor\)。 (R/W)
\end{reglist}\end{regdesc}
\end{register}


\begin{register}{H}{ASRC\_CHNL1\_RESET\_REG}{0x0048}\label{regdesc:ASRCCHNL1RESETREG}
\regfield{(reserved)}{31}{1}{0000000000000000000000000000000}%
\regfield{ASRC\_CHNL1\_RESET}{1}{0}{{0}}%
\reglabel{Reset}\regnewline%
\begin{regdesc}\begin{reglist}
\label{fielddesc:ASRCCHNL1RESET}\item [ASRC\_CHNL1\_RESET] 设置此位以复位通道 1。 (WT)
\end{reglist}\end{regdesc}
\end{register}


\begin{register}{H}{ASRC\_CHNL1\_START\_REG}{0x004C}\label{regdesc:ASRCCHNL1STARTREG}
\regfield{(reserved)}{31}{1}{0000000000000000000000000000000}%
\regfield{ASRC\_CHNL1\_START}{1}{0}{{0}}%
\reglabel{Reset}\regnewline%
\begin{regdesc}\begin{reglist}
\label{fielddesc:ASRCCHNL1START}\item [ASRC\_CHNL1\_START] 设置此位以启动通道 1。 (R/W)
\end{reglist}\end{regdesc}
\end{register}


\begin{register}{H}{ASRC\_CHNL1\_DIN\_CFG\_REG}{0x0050}\label{regdesc:ASRCCHNL1DINCFGREG}
\regfield{ASRC\_CHNL1\_IN\_LEN}{24}{8}{{0}}%
\regfield{(reserved)}{6}{2}{000000}%
\regfield{ASRC\_CHNL1\_IN\_CNT\_CLR}{1}{1}{{0}}%
\regfield{ASRC\_CHNL1\_IN\_CNT\_ENA}{1}{0}{{0}}%
\reglabel{Reset}\regnewline%
\begin{regdesc}\begin{reglist}
\label{fielddesc:ASRCCHNL1INCNTENA}\item [ASRC\_CHNL1\_IN\_CNT\_ENA] 设置此位以使能输入数据字节计数器。 (R/W)
\label{fielddesc:ASRCCHNL1INCNTCLR}\item [ASRC\_CHNL1\_IN\_CNT\_CLR] 设置此位以清除输入数据字节计数器。 (WT)
\label{fielddesc:ASRCCHNL1INLEN}\item [ASRC\_CHNL1\_IN\_LEN] 写入该字段以指定从 DMA 接收的数据字节数。 (R/W)
\end{reglist}\end{regdesc}
\end{register}


\begin{register}{H}{ASRC\_CHNL1\_DOUT\_CFG\_REG}{0x0054}\label{regdesc:ASRCCHNL1DOUTCFGREG}
\regfield{ASRC\_CHNL1\_OUT\_LEN}{24}{8}{{0}}%
\regfield{(reserved)}{3}{5}{000}%
\regfield{ASRC\_CHNL1\_OUT\_SAMPLES\_COMP}{1}{4}{{0}}%
\regfield{ASRC\_CHNL1\_OUT\_CNT\_CLR}{1}{3}{{0}}%
\regfield{ASRC\_CHNL1\_OUT\_CNT\_ENA}{1}{2}{{0}}%
\regfield{ASRC\_CHNL1\_OUT\_EOF\_GEN\_MODE}{2}{0}{{0}}%
\reglabel{Reset}\regnewline%
\begin{regdesc}\begin{reglist}
\label{fielddesc:ASRCCHNL1OUTEOFGENMODE}\item [ASRC\_CHNL1\_OUT\_EOF\_GEN\_MODE] 写入该字段以配置写入 DMA 的 EOF 类型：\\
0：计数器 EOF 事件\\
1：DMA 输入 EOF 信号\\
2：计数器 EOF 和 DMA 输入 EOF 均写入\\
3：无 EOF 写入\\
(R/W)
\label{fielddesc:ASRCCHNL1OUTCNTENA}\item [ASRC\_CHNL1\_OUT\_CNT\_ENA] 设置此位以启用输出数据字节计数器。 (R/W)
\label{fielddesc:ASRCCHNL1OUTCNTCLR}\item [ASRC\_CHNL1\_OUT\_CNT\_CLR] 写入此位以清除输出数据字节计数器。 (WT)
\label{fielddesc:ASRCCHNL1OUTSAMPLESCOMP}\item [ASRC\_CHNL1\_OUT\_SAMPLES\_COMP] 设置此位以启用输出数据字节计数器补偿，当使用分数重采样器并进行 2 倍抽取且 ASRC\_CHNL1\_OUT\_CNT ≥ ASRC\_CHNL1\_OUT\_LEN 时生效。 (R/W)
\label{fielddesc:ASRCCHNL1OUTLEN}\item [ASRC\_CHNL1\_OUT\_LEN] 写入该字段以指定传输到 DMA 的数据字节数，当计数器到达末尾时，EOF（结束标志）将被置位。 (R/W)
\end{reglist}\end{regdesc}
\end{register}



\begin{register}{H}{ASRC\_CHNL1\_FIFO\_CTRL\_REG}{0x0058}\label{regdesc:ASRCCHNL1FIFOCTRLREG}
\regfield{(reserved)}{30}{2}{000000000000000000000000000000}%
\regfield{ASRC\_CHNL1\_OUTFIFO\_RESET}{1}{1}{{0}}%
\regfield{ASRC\_CHNL1\_INFIFO\_RESET}{1}{0}{{0}}%
\reglabel{Reset}\regnewline%
\begin{regdesc}\begin{reglist}
\label{fielddesc:ASRCCHNL1INFIFORESET}\item [ASRC\_CHNL1\_INFIFO\_RESET] 写入此位以复位输入 FIFO。 (WT)
\label{fielddesc:ASRCCHNL1OUTFIFORESET}\item [ASRC\_CHNL1\_OUTFIFO\_RESET] 写入此位以复位输出 FIFO。 (WT)


\end{reglist}\end{regdesc}
\end{register}


\begin{register}{H}{ASRC\_CHNL0\_INT\_RAW\_REG}{0x0020}\label{regdesc:ASRCCHNL0INTRAWREG}
\regfield{(reserved)}{31}{1}{0000000000000000000000000000000}%
\regfield{ASRC\_CHNL0\_OUTCNT\_EOF\_INT\_RAW}{1}{0}{{0}}%
\reglabel{Reset}\regnewline%
\begin{regdesc}\begin{reglist}
\label{fielddesc:ASRCCHNL0OUTCNTEOFINTRAW}\item [ASRC\_CHNL0\_OUTCNT\_EOF\_INT\_RAW] \hyperref[int:ASRCCHNL0INT] {ASRC\_CHNL0\_INT} 中断的原始中断状态。(R/WTC/SS)
\end{reglist}\end{regdesc}
\end{register}


\begin{register}{H}{ASRC\_CHNL0\_INT\_ST\_REG}{0x0024}\label{regdesc:ASRCCHNL0INTSTREG}
\regfield{(reserved)}{31}{1}{0000000000000000000000000000000}%
\regfield{ASRC\_CHNL0\_OUTCNT\_EOF\_INT\_ST}{1}{0}{{0}}%
\reglabel{Reset}\regnewline%
\begin{regdesc}\begin{reglist}
\label{fielddesc:ASRCCHNL0OUTCNTEOFINTST}\item [ASRC\_CHNL0\_OUTCNT\_EOF\_INT\_ST] \hyperref[int:ASRCCHNL0INT] {ASRC\_CHNL0\_INT} 中断的屏蔽中断状态。 (RO)
\end{reglist}\end{regdesc}
\end{register}


\begin{register}{H}{ASRC\_CHNL0\_INT\_ENA\_REG}{0x0028}\label{regdesc:ASRCCHNL0INTENAREG}
\regfield{(reserved)}{31}{1}{0000000000000000000000000000000}%
\regfield{ASRC\_CHNL0\_OUTCNT\_EOF\_INT\_ENA}{1}{0}{{0}}%
\reglabel{Reset}\regnewline%
\begin{regdesc}\begin{reglist}
\label{fielddesc:ASRCCHNL0OUTCNTEOFINTENA}\item [ASRC\_CHNL0\_OUTCNT\_EOF\_INT\_ENA] 写 1 使能 \hyperref[int:ASRCCHNL0INT] {ASRC\_CHNL0\_INT} 中断。 (R/W)
\end{reglist}\end{regdesc}
\end{register}


\begin{register}{H}{ASRC\_CHNL0\_INT\_CLR\_REG}{0x002C}\label{regdesc:ASRCCHNL0INTCLRREG}
\regfield{(reserved)}{31}{1}{0000000000000000000000000000000}%
\regfield{ASRC\_CHNL0\_OUTCNT\_EOF\_INT\_CLR}{1}{0}{{0}}%
\reglabel{Reset}\regnewline%
\begin{regdesc}\begin{reglist}
\label{fielddesc:ASRCCHNL0OUTCNTEOFINTCLR}\item [ASRC\_CHNL0\_OUTCNT\_EOF\_INT\_CLR] 写 1 清除 \hyperref[int:ASRCCHNL0INT] {ASRC\_CHNL0\_INT} 中断。 (WT)
\end{reglist}\end{regdesc}
\end{register}


\begin{register}{H}{ASRC\_CHNL1\_INT\_RAW\_REG}{0x005C}\label{regdesc:ASRCCHNL1INTRAWREG}
\regfield{(reserved)}{31}{1}{0000000000000000000000000000000}%
\regfield{ASRC\_CHNL1\_OUTCNT\_EOF\_INT\_RAW}{1}{0}{{0}}%
\reglabel{Reset}\regnewline%
\begin{regdesc}\begin{reglist}
\label{fielddesc:ASRCCHNL1OUTCNTEOFINTRAW}\item [ASRC\_CHNL1\_OUTCNT\_EOF\_INT\_RAW] \hyperref[int:ASRCCHNL1INT] {ASRC\_CHNL1\_INT} 中断的原始中断状态。 (R/WTC/SS)
\end{reglist}\end{regdesc}
\end{register}


\begin{register}{H}{ASRC\_CHNL1\_INT\_ST\_REG}{0x0060}\label{regdesc:ASRCCHNL1INTSTREG}
\regfield{(reserved)}{31}{1}{0000000000000000000000000000000}%
\regfield{ASRC\_CHNL1\_OUTCNT\_EOF\_INT\_ST}{1}{0}{{0}}%
\reglabel{Reset}\regnewline%
\begin{regdesc}\begin{reglist}
\label{fielddesc:ASRCCHNL1OUTCNTEOFINTST}\item [ASRC\_CHNL1\_OUTCNT\_EOF\_INT\_ST] \hyperref[int:ASRCCHNL1INT] {ASRC\_CHNL1\_INT} 中断的屏蔽中断状态。 (RO)
\end{reglist}\end{regdesc}
\end{register}


\begin{register}{H}{ASRC\_CHNL1\_INT\_ENA\_REG}{0x0064}\label{regdesc:ASRCCHNL1INTENAREG}
\regfield{(reserved)}{31}{1}{0000000000000000000000000000000}%
\regfield{ASRC\_CHNL1\_OUTCNT\_EOF\_INT\_ENA}{1}{0}{{0}}%
\reglabel{Reset}\regnewline%
\begin{regdesc}\begin{reglist}
\label{fielddesc:ASRCCHNL1OUTCNTEOFINTENA}\item [ASRC\_CHNL1\_OUTCNT\_EOF\_INT\_ENA] 写 1 使能 \hyperref[int:ASRCCHNL1INT] {ASRC\_CHNL1\_INT} 中断。 (R/W)
\end{reglist}\end{regdesc}
\end{register}


\begin{register}{H}{ASRC\_CHNL1\_INT\_CLR\_REG}{0x0068}\label{regdesc:ASRCCHNL1INTCLRREG}
\regfield{(reserved)}{31}{1}{0000000000000000000000000000000}%
\regfield{ASRC\_CHNL1\_OUTCNT\_EOF\_INT\_CLR}{1}{0}{{0}}%
\reglabel{Reset}\regnewline%
\begin{regdesc}\begin{reglist}
\label{fielddesc:ASRCCHNL1OUTCNTEOFINTCLR}\item [ASRC\_CHNL1\_OUTCNT\_EOF\_INT\_CLR] 写 1 清除 \hyperref[int:ASRCCHNL1INT] {ASRC\_CHNL1\_INT} 中断。 (WT)
\end{reglist}\end{regdesc}
\end{register}


\begin{register}{H}{ASRC\_CHNL0\_OUT\_CNT\_REG}{0x0030}\label{regdesc:ASRCCHNL0OUTCNTREG}
\regfield{(reserved)}{8}{24}{00000000}%
\regfield{ASRC\_CHNL0\_OUT\_CNT}{24}{0}{{0}}%
\reglabel{Reset}\regnewline%
\begin{regdesc}\begin{reglist}
\label{fielddesc:ASRCCHNL0OUTCNT}\item [ASRC\_CHNL0\_OUT\_CNT] 累计向 DMA 写入的总字节数。 (RO)
\end{reglist}\end{regdesc}
\end{register}


\begin{register}{H}{ASRC\_CHNL1\_OUT\_CNT\_REG}{0x006C}\label{regdesc:ASRCCHNL1OUTCNTREG}
\regfield{(reserved)}{8}{24}{00000000}%
\regfield{ASRC\_CHNL1\_OUT\_CNT}{24}{0}{{0}}%
\reglabel{Reset}\regnewline%
\begin{regdesc}\begin{reglist}
\label{fielddesc:ASRCCHNL1OUTCNT}\item [ASRC\_CHNL1\_OUT\_CNT] 累计向 DMA 写入的总字节数。 (RO)
\end{reglist}\end{regdesc}
\end{register}


\begin{register}{H}{ASRC\_CHNL0\_TRACE1\_REG}{0x0038}\label{regdesc:ASRCCHNL0TRACE1REG}
\regfield{(reserved)}{8}{24}{00000000}%
\regfield{ASRC\_CHNL0\_OUT\_INC}{24}{0}{{0}}%
\reglabel{Reset}\regnewline%
\begin{regdesc}\begin{reglist}
\label{fielddesc:ASRCCHNL0OUTINC}\item [ASRC\_CHNL0\_OUT\_INC] 自上次 EOF 发生后向 DMA 写入的字节数。 (RO)
\end{reglist}\end{regdesc}
\end{register}


\begin{register}{H}{ASRC\_CHNL1\_TRACE1\_REG}{0x0074}\label{regdesc:ASRCCHNL1TRACE1REG}
\regfield{(reserved)}{8}{24}{00000000}%
\regfield{ASRC\_CHNL1\_OUT\_INC}{24}{0}{{0}}%
\reglabel{Reset}\regnewline%
\begin{regdesc}\begin{reglist}
\label{fielddesc:ASRCCHNL1OUTINC}\item [ASRC\_CHNL1\_OUT\_INC] 自上次 EOF 发生后向 DMA 写入的字节数。 (RO)
\end{reglist}\end{regdesc}
\end{register}


\begin{register}{H}{ASRC\_SYS\_REG}{0x00F8}\label{regdesc:ASRCSYSREG}
\regfield{(reserved)}{25}{7}{0000000000000000000000000}%
\regfield{ASRC\_CHNL1\_INFIFO\_CLK\_FO}{1}{6}{{0}}%
\regfield{ASRC\_CHNL1\_OUTFIFO\_CLK\_FO}{1}{5}{{0}}%
\regfield{ASRC\_CHNL0\_INFIFO\_CLK\_FO}{1}{4}{{0}}%
\regfield{ASRC\_CHNL0\_OUTFIFO\_CLK\_FO}{1}{3}{{0}}%
\regfield{ASRC\_CHNL1\_CLK\_FO}{1}{2}{{0}}%
\regfield{ASRC\_CHNL0\_CLK\_FO}{1}{1}{{0}}%
\regfield{ASRC\_CLK\_EN}{1}{0}{{0}}%
\reglabel{Reset}\regnewline%
\begin{regdesc}\begin{reglist}
\label{fielddesc:ASRCCLKEN}\item [ASRC\_CLK\_EN] 配置是否启用 ASRC 模块的时钟。\\
0: 禁用 \\
1: 启用 \\
(R/W)
\label{fielddesc:ASRCCHNL0CLKFO}\item [ASRC\_CHNL0\_CLK\_FO] 设置此位使通道0时钟自由运行。 (R/W)
\label{fielddesc:ASRCCHNL1CLKFO}\item [ASRC\_CHNL1\_CLK\_FO] 设置此位使通道1时钟自由运行。 (R/W)
\label{fielddesc:ASRCCHNL0OUTFIFOCLKFO}\item [ASRC\_CHNL0\_OUTFIFO\_CLK\_FO] 设置此位使通道0输出 FIFO 时钟自由运行。 (R/W)
\label{fielddesc:ASRCCHNL0INFIFOCLKFO}\item [ASRC\_CHNL0\_INFIFO\_CLK\_FO] 设置此位使通道0输入 FIFO 时钟自由运行。 (R/W)
\label{fielddesc:ASRCCHNL1OUTFIFOCLKFO}\item [ASRC\_CHNL1\_OUTFIFO\_CLK\_FO] 设置此位使通道1输出 FIFO 时钟自由运行。 (R/W)
\label{fielddesc:ASRCCHNL1INFIFOCLKFO}\item [ASRC\_CHNL1\_INFIFO\_CLK\_FO] 设置此位使通道1输入 FIFO 时钟自由运行。 (R/W)
\end{reglist}\end{regdesc}
\end{register}


\begin{register}{H}{ASRC\_DATE\_REG}{0x00FC}\label{regdesc:ASRCDATEREG}
\regfield{(reserved)}{4}{28}{0000}%
\regfield{ASRC\_DATE}{28}{0}{{0x2407240}}%
\reglabel{Reset}\regnewline%
\begin{regdesc}\begin{reglist}
\label{fielddesc:ASRCDATE}\item [ASRC\_DATE] 版本控制寄存器。 (R/W)
\end{reglist}\end{regdesc}
\end{register}
